% =====================================================================
% Pruebas funcionales de RF-01 (Periodos y categorias documentales) - Roselyn
% FRAGMENTO para el informe: % =====================================================================
% Pruebas funcionales de RF-01 (Periodos y categorias documentales) - Roselyn
% FRAGMENTO para el informe: % =====================================================================
% Pruebas funcionales de RF-01 (Periodos y categorias documentales) - Roselyn
% FRAGMENTO para el informe: % =====================================================================
% Pruebas funcionales de RF-01 (Periodos y categorias documentales) - Roselyn
% FRAGMENTO para el informe: \input{rf01_tabla_pruebas}
% Paquetes: booktabs, longtable, array. Las capturas estan en la carpeta evidencias/.
% =====================================================================

\providecommand{\evus}{\textunderscore\allowbreak}
\providecommand{\ev}[1]{{\texttt{\scriptsize\let\_\evus #1}}}

\subsection{Pruebas funcionales de RF-01: periodos y categorías documentales}

Cada prueba se hizo desde la aplicación y después se comprobó en PostgreSQL con una
consulta, para confirmar que lo que se vio en pantalla coincide con lo guardado en la base
de datos. Las capturas están en la carpeta \texttt{evidencias/} del repositorio.

\begingroup\footnotesize
\setlength{\tabcolsep}{3pt}
\begin{longtable}{>{\raggedright\arraybackslash}p{1.5cm}>{\raggedright\arraybackslash}p{1.8cm}>{\raggedright\arraybackslash}p{2.8cm}>{\raggedright\arraybackslash}p{2.7cm}>{\raggedright\arraybackslash}p{2.7cm}>{\raggedright\arraybackslash}p{3.2cm}}
\caption{Pruebas funcionales de RF-01}\label{tab:pruebas-rf01}\\
\toprule
\textbf{Requisito} & \textbf{Prueba} & \textbf{Datos usados} & \textbf{Resultado esperado} & \textbf{Resultado obtenido} & \textbf{Evidencia} \\
\midrule
\endfirsthead
\toprule
\textbf{Requisito} & \textbf{Prueba} & \textbf{Datos usados} & \textbf{Resultado esperado} & \textbf{Resultado obtenido} & \textbf{Evidencia} \\
\midrule
\endhead
RF-01 & Crear categoría &
Nombre: Plan de acción inicial; activa. &
Se guarda y aparece en la lista con el mensaje ``Categoria creada''. &
Apareció el mensaje y la categoría en la lista. &
\ev{RF01\_categoria\_crear.png} \\
\addlinespace
RF-01 & Categoría repetida &
Nombre igual al de una categoría existente. &
Mensaje de error y no se guarda. &
Mostró ``Ya existe una categoria con ese nombre'' y no guardó. &
\ev{RF01\_categoria\_nombre\_repetido\_error.png} \\
\addlinespace
RF-01 & Categoría sin nombre &
Nombre vacío. &
Mensaje de error y no se guarda. &
Mostró ``El nombre es obligatorio'' y no guardó. &
\ev{RF01\_categoria\_nombre\_vacio\_error.png} \\
\addlinespace
RF-01 & Editar categoría &
Cambio de la descripción de una categoría. &
Se actualiza con el mensaje ``Categoria actualizada''. &
Se actualizó y apareció el mensaje. &
\ev{RF01\_categoria\_editar.png} \\
\addlinespace
RF-01 & Eliminar categoría &
Categoría que no está en uso. &
Se elimina con el mensaje ``Categoria eliminada''. &
Se eliminó y apareció el mensaje. &
\ev{RF01\_categoria\_eliminar.png} \\
\addlinespace
RF-01 & Crear periodo &
Nombre: 2026-2027 CI; fechas válidas; tres categorías marcadas. &
Se guarda con sus categorías y aparece en la lista con el mensaje ``Periodo creado''. &
Apareció el mensaje y el periodo en la lista con tres categorías. &
\ev{RF01\_crear\_periodo.png} \\
\addlinespace
RF-01 & Reutilizar categorías &
Periodo nuevo con las categorías del periodo anterior marcadas y una desmarcada. &
Se guarda solo con las categorías marcadas. &
Se guardó con las categorías elegidas. En PostgreSQL, \texttt{periodo\_categoria} muestra solo esas asignaciones. &
\ev{RF01\_reutilizar\_categorias.png}\newline
\ev{RF01\_reutilizar\_categorias\_lista.png}\newline
\ev{RF01\_reutilizar\_categorias\_postgres.png} \\
\addlinespace
RF-01 & Periodo con categoría nueva &
Periodo nuevo y una categoría nueva escrita en el campo de categorías nuevas. &
Se crea la categoría y queda asignada al periodo. &
La categoría nueva se creó y quedó asignada al periodo. &
\ev{RF01\_periodo\_crear\_con\_categoria\_nueva.png} \\
\addlinespace
RF-01 & Fechas inválidas &
Fecha de fin anterior a la fecha de inicio. &
Mensaje de error y no se guarda. &
Mostró el mensaje de fecha de fin posterior a la de inicio y no guardó. &
\ev{RF01\_fechas\_error.png} \\
\addlinespace
RF-01 & Periodo repetido &
Nombre de un periodo que ya existe. &
Mensaje de error y no se guarda. &
Mostró ``Ya existe un periodo con ese nombre'' y no guardó. &
\ev{RF01\_periodo\_nombre\_repetido\_error.png} \\
\addlinespace
RF-01 & Editar periodo &
Cambio de nombre o de fechas. &
Se actualiza con el mensaje ``Periodo actualizado''. &
Se actualizó y apareció el mensaje. &
\ev{RF01\_periodo\_editar.png} \\
\addlinespace
RF-01 & Cerrar periodo &
Periodo en estado ACTIVO. &
El estado cambia a CERRADO. &
El estado cambió a CERRADO. &
\ev{RF01\_periodo\_cerrar.png} \\
\addlinespace
RF-01 & Eliminar periodo &
Periodo de prueba. &
Se elimina con el mensaje ``Periodo eliminado''. &
Se eliminó y apareció el mensaje. &
\ev{RF01\_periodo\_eliminar.png} \\
\addlinespace
RF-01 & Verificación en PostgreSQL &
Consulta con \texttt{JOIN} entre \texttt{periodo}, \texttt{periodo\_categoria} y \texttt{categoria\_documental}. &
Los datos de la base coinciden con lo hecho en la aplicación. &
Los registros coincidieron con lo hecho en la aplicación. &
\ev{RF01\_verificacion\_postgres\_periodo\_categoria.png} \\
\bottomrule
\end{longtable}
\endgroup
% Paquetes: booktabs, longtable, array. Las capturas estan en la carpeta evidencias/.
% =====================================================================

\providecommand{\evus}{\textunderscore\allowbreak}
\providecommand{\ev}[1]{{\texttt{\scriptsize\let\_\evus #1}}}

\subsection{Pruebas funcionales de RF-01: periodos y categorías documentales}

Cada prueba se hizo desde la aplicación y después se comprobó en PostgreSQL con una
consulta, para confirmar que lo que se vio en pantalla coincide con lo guardado en la base
de datos. Las capturas están en la carpeta \texttt{evidencias/} del repositorio.

\begingroup\footnotesize
\setlength{\tabcolsep}{3pt}
\begin{longtable}{>{\raggedright\arraybackslash}p{1.5cm}>{\raggedright\arraybackslash}p{1.8cm}>{\raggedright\arraybackslash}p{2.8cm}>{\raggedright\arraybackslash}p{2.7cm}>{\raggedright\arraybackslash}p{2.7cm}>{\raggedright\arraybackslash}p{3.2cm}}
\caption{Pruebas funcionales de RF-01}\label{tab:pruebas-rf01}\\
\toprule
\textbf{Requisito} & \textbf{Prueba} & \textbf{Datos usados} & \textbf{Resultado esperado} & \textbf{Resultado obtenido} & \textbf{Evidencia} \\
\midrule
\endfirsthead
\toprule
\textbf{Requisito} & \textbf{Prueba} & \textbf{Datos usados} & \textbf{Resultado esperado} & \textbf{Resultado obtenido} & \textbf{Evidencia} \\
\midrule
\endhead
RF-01 & Crear categoría &
Nombre: Plan de acción inicial; activa. &
Se guarda y aparece en la lista con el mensaje ``Categoria creada''. &
Apareció el mensaje y la categoría en la lista. &
\ev{RF01\_categoria\_crear.png} \\
\addlinespace
RF-01 & Categoría repetida &
Nombre igual al de una categoría existente. &
Mensaje de error y no se guarda. &
Mostró ``Ya existe una categoria con ese nombre'' y no guardó. &
\ev{RF01\_categoria\_nombre\_repetido\_error.png} \\
\addlinespace
RF-01 & Categoría sin nombre &
Nombre vacío. &
Mensaje de error y no se guarda. &
Mostró ``El nombre es obligatorio'' y no guardó. &
\ev{RF01\_categoria\_nombre\_vacio\_error.png} \\
\addlinespace
RF-01 & Editar categoría &
Cambio de la descripción de una categoría. &
Se actualiza con el mensaje ``Categoria actualizada''. &
Se actualizó y apareció el mensaje. &
\ev{RF01\_categoria\_editar.png} \\
\addlinespace
RF-01 & Eliminar categoría &
Categoría que no está en uso. &
Se elimina con el mensaje ``Categoria eliminada''. &
Se eliminó y apareció el mensaje. &
\ev{RF01\_categoria\_eliminar.png} \\
\addlinespace
RF-01 & Crear periodo &
Nombre: 2026-2027 CI; fechas válidas; tres categorías marcadas. &
Se guarda con sus categorías y aparece en la lista con el mensaje ``Periodo creado''. &
Apareció el mensaje y el periodo en la lista con tres categorías. &
\ev{RF01\_crear\_periodo.png} \\
\addlinespace
RF-01 & Reutilizar categorías &
Periodo nuevo con las categorías del periodo anterior marcadas y una desmarcada. &
Se guarda solo con las categorías marcadas. &
Se guardó con las categorías elegidas. En PostgreSQL, \texttt{periodo\_categoria} muestra solo esas asignaciones. &
\ev{RF01\_reutilizar\_categorias.png}\newline
\ev{RF01\_reutilizar\_categorias\_lista.png}\newline
\ev{RF01\_reutilizar\_categorias\_postgres.png} \\
\addlinespace
RF-01 & Periodo con categoría nueva &
Periodo nuevo y una categoría nueva escrita en el campo de categorías nuevas. &
Se crea la categoría y queda asignada al periodo. &
La categoría nueva se creó y quedó asignada al periodo. &
\ev{RF01\_periodo\_crear\_con\_categoria\_nueva.png} \\
\addlinespace
RF-01 & Fechas inválidas &
Fecha de fin anterior a la fecha de inicio. &
Mensaje de error y no se guarda. &
Mostró el mensaje de fecha de fin posterior a la de inicio y no guardó. &
\ev{RF01\_fechas\_error.png} \\
\addlinespace
RF-01 & Periodo repetido &
Nombre de un periodo que ya existe. &
Mensaje de error y no se guarda. &
Mostró ``Ya existe un periodo con ese nombre'' y no guardó. &
\ev{RF01\_periodo\_nombre\_repetido\_error.png} \\
\addlinespace
RF-01 & Editar periodo &
Cambio de nombre o de fechas. &
Se actualiza con el mensaje ``Periodo actualizado''. &
Se actualizó y apareció el mensaje. &
\ev{RF01\_periodo\_editar.png} \\
\addlinespace
RF-01 & Cerrar periodo &
Periodo en estado ACTIVO. &
El estado cambia a CERRADO. &
El estado cambió a CERRADO. &
\ev{RF01\_periodo\_cerrar.png} \\
\addlinespace
RF-01 & Eliminar periodo &
Periodo de prueba. &
Se elimina con el mensaje ``Periodo eliminado''. &
Se eliminó y apareció el mensaje. &
\ev{RF01\_periodo\_eliminar.png} \\
\addlinespace
RF-01 & Verificación en PostgreSQL &
Consulta con \texttt{JOIN} entre \texttt{periodo}, \texttt{periodo\_categoria} y \texttt{categoria\_documental}. &
Los datos de la base coinciden con lo hecho en la aplicación. &
Los registros coincidieron con lo hecho en la aplicación. &
\ev{RF01\_verificacion\_postgres\_periodo\_categoria.png} \\
\bottomrule
\end{longtable}
\endgroup
% Paquetes: booktabs, longtable, array. Las capturas estan en la carpeta evidencias/.
% =====================================================================

\providecommand{\evus}{\textunderscore\allowbreak}
\providecommand{\ev}[1]{{\texttt{\scriptsize\let\_\evus #1}}}

\subsection{Pruebas funcionales de RF-01: periodos y categorías documentales}

Cada prueba se hizo desde la aplicación y después se comprobó en PostgreSQL con una
consulta, para confirmar que lo que se vio en pantalla coincide con lo guardado en la base
de datos. Las capturas están en la carpeta \texttt{evidencias/} del repositorio.

\begingroup\footnotesize
\setlength{\tabcolsep}{3pt}
\begin{longtable}{>{\raggedright\arraybackslash}p{1.5cm}>{\raggedright\arraybackslash}p{1.8cm}>{\raggedright\arraybackslash}p{2.8cm}>{\raggedright\arraybackslash}p{2.7cm}>{\raggedright\arraybackslash}p{2.7cm}>{\raggedright\arraybackslash}p{3.2cm}}
\caption{Pruebas funcionales de RF-01}\label{tab:pruebas-rf01}\\
\toprule
\textbf{Requisito} & \textbf{Prueba} & \textbf{Datos usados} & \textbf{Resultado esperado} & \textbf{Resultado obtenido} & \textbf{Evidencia} \\
\midrule
\endfirsthead
\toprule
\textbf{Requisito} & \textbf{Prueba} & \textbf{Datos usados} & \textbf{Resultado esperado} & \textbf{Resultado obtenido} & \textbf{Evidencia} \\
\midrule
\endhead
RF-01 & Crear categoría &
Nombre: Plan de acción inicial; activa. &
Se guarda y aparece en la lista con el mensaje ``Categoria creada''. &
Apareció el mensaje y la categoría en la lista. &
\ev{RF01\_categoria\_crear.png} \\
\addlinespace
RF-01 & Categoría repetida &
Nombre igual al de una categoría existente. &
Mensaje de error y no se guarda. &
Mostró ``Ya existe una categoria con ese nombre'' y no guardó. &
\ev{RF01\_categoria\_nombre\_repetido\_error.png} \\
\addlinespace
RF-01 & Categoría sin nombre &
Nombre vacío. &
Mensaje de error y no se guarda. &
Mostró ``El nombre es obligatorio'' y no guardó. &
\ev{RF01\_categoria\_nombre\_vacio\_error.png} \\
\addlinespace
RF-01 & Editar categoría &
Cambio de la descripción de una categoría. &
Se actualiza con el mensaje ``Categoria actualizada''. &
Se actualizó y apareció el mensaje. &
\ev{RF01\_categoria\_editar.png} \\
\addlinespace
RF-01 & Eliminar categoría &
Categoría que no está en uso. &
Se elimina con el mensaje ``Categoria eliminada''. &
Se eliminó y apareció el mensaje. &
\ev{RF01\_categoria\_eliminar.png} \\
\addlinespace
RF-01 & Crear periodo &
Nombre: 2026-2027 CI; fechas válidas; tres categorías marcadas. &
Se guarda con sus categorías y aparece en la lista con el mensaje ``Periodo creado''. &
Apareció el mensaje y el periodo en la lista con tres categorías. &
\ev{RF01\_crear\_periodo.png} \\
\addlinespace
RF-01 & Reutilizar categorías &
Periodo nuevo con las categorías del periodo anterior marcadas y una desmarcada. &
Se guarda solo con las categorías marcadas. &
Se guardó con las categorías elegidas. En PostgreSQL, \texttt{periodo\_categoria} muestra solo esas asignaciones. &
\ev{RF01\_reutilizar\_categorias.png}\newline
\ev{RF01\_reutilizar\_categorias\_lista.png}\newline
\ev{RF01\_reutilizar\_categorias\_postgres.png} \\
\addlinespace
RF-01 & Periodo con categoría nueva &
Periodo nuevo y una categoría nueva escrita en el campo de categorías nuevas. &
Se crea la categoría y queda asignada al periodo. &
La categoría nueva se creó y quedó asignada al periodo. &
\ev{RF01\_periodo\_crear\_con\_categoria\_nueva.png} \\
\addlinespace
RF-01 & Fechas inválidas &
Fecha de fin anterior a la fecha de inicio. &
Mensaje de error y no se guarda. &
Mostró el mensaje de fecha de fin posterior a la de inicio y no guardó. &
\ev{RF01\_fechas\_error.png} \\
\addlinespace
RF-01 & Periodo repetido &
Nombre de un periodo que ya existe. &
Mensaje de error y no se guarda. &
Mostró ``Ya existe un periodo con ese nombre'' y no guardó. &
\ev{RF01\_periodo\_nombre\_repetido\_error.png} \\
\addlinespace
RF-01 & Editar periodo &
Cambio de nombre o de fechas. &
Se actualiza con el mensaje ``Periodo actualizado''. &
Se actualizó y apareció el mensaje. &
\ev{RF01\_periodo\_editar.png} \\
\addlinespace
RF-01 & Cerrar periodo &
Periodo en estado ACTIVO. &
El estado cambia a CERRADO. &
El estado cambió a CERRADO. &
\ev{RF01\_periodo\_cerrar.png} \\
\addlinespace
RF-01 & Eliminar periodo &
Periodo de prueba. &
Se elimina con el mensaje ``Periodo eliminado''. &
Se eliminó y apareció el mensaje. &
\ev{RF01\_periodo\_eliminar.png} \\
\addlinespace
RF-01 & Verificación en PostgreSQL &
Consulta con \texttt{JOIN} entre \texttt{periodo}, \texttt{periodo\_categoria} y \texttt{categoria\_documental}. &
Los datos de la base coinciden con lo hecho en la aplicación. &
Los registros coincidieron con lo hecho en la aplicación. &
\ev{RF01\_verificacion\_postgres\_periodo\_categoria.png} \\
\bottomrule
\end{longtable}
\endgroup
% Paquetes: booktabs, longtable, array. Las capturas estan en la carpeta evidencias/.
% =====================================================================

\providecommand{\evus}{\textunderscore\allowbreak}
\providecommand{\ev}[1]{{\texttt{\scriptsize\let\_\evus #1}}}

\subsection{Pruebas funcionales de RF-01: periodos y categorías documentales}

Cada prueba se hizo desde la aplicación y después se comprobó en PostgreSQL con una
consulta, para confirmar que lo que se vio en pantalla coincide con lo guardado en la base
de datos. Las capturas están en la carpeta \texttt{evidencias/} del repositorio.

\begingroup\footnotesize
\setlength{\tabcolsep}{3pt}
\begin{longtable}{>{\raggedright\arraybackslash}p{1.5cm}>{\raggedright\arraybackslash}p{1.8cm}>{\raggedright\arraybackslash}p{2.8cm}>{\raggedright\arraybackslash}p{2.7cm}>{\raggedright\arraybackslash}p{2.7cm}>{\raggedright\arraybackslash}p{3.2cm}}
\caption{Pruebas funcionales de RF-01}\label{tab:pruebas-rf01}\\
\toprule
\textbf{Requisito} & \textbf{Prueba} & \textbf{Datos usados} & \textbf{Resultado esperado} & \textbf{Resultado obtenido} & \textbf{Evidencia} \\
\midrule
\endfirsthead
\toprule
\textbf{Requisito} & \textbf{Prueba} & \textbf{Datos usados} & \textbf{Resultado esperado} & \textbf{Resultado obtenido} & \textbf{Evidencia} \\
\midrule
\endhead
RF-01 & Crear categoría &
Nombre: Plan de acción inicial; activa. &
Se guarda y aparece en la lista con el mensaje ``Categoria creada''. &
Apareció el mensaje y la categoría en la lista. &
\ev{RF01\_categoria\_crear.png} \\
\addlinespace
RF-01 & Categoría repetida &
Nombre igual al de una categoría existente. &
Mensaje de error y no se guarda. &
Mostró ``Ya existe una categoria con ese nombre'' y no guardó. &
\ev{RF01\_categoria\_nombre\_repetido\_error.png} \\
\addlinespace
RF-01 & Categoría sin nombre &
Nombre vacío. &
Mensaje de error y no se guarda. &
Mostró ``El nombre es obligatorio'' y no guardó. &
\ev{RF01\_categoria\_nombre\_vacio\_error.png} \\
\addlinespace
RF-01 & Editar categoría &
Cambio de la descripción de una categoría. &
Se actualiza con el mensaje ``Categoria actualizada''. &
Se actualizó y apareció el mensaje. &
\ev{RF01\_categoria\_editar.png} \\
\addlinespace
RF-01 & Eliminar categoría &
Categoría que no está en uso. &
Se elimina con el mensaje ``Categoria eliminada''. &
Se eliminó y apareció el mensaje. &
\ev{RF01\_categoria\_eliminar.png} \\
\addlinespace
RF-01 & Crear periodo &
Nombre: 2026-2027 CI; fechas válidas; tres categorías marcadas. &
Se guarda con sus categorías y aparece en la lista con el mensaje ``Periodo creado''. &
Apareció el mensaje y el periodo en la lista con tres categorías. &
\ev{RF01\_crear\_periodo.png} \\
\addlinespace
RF-01 & Reutilizar categorías &
Periodo nuevo con las categorías del periodo anterior marcadas y una desmarcada. &
Se guarda solo con las categorías marcadas. &
Se guardó con las categorías elegidas. En PostgreSQL, \texttt{periodo\_categoria} muestra solo esas asignaciones. &
\ev{RF01\_reutilizar\_categorias.png}\newline
\ev{RF01\_reutilizar\_categorias\_lista.png}\newline
\ev{RF01\_reutilizar\_categorias\_postgres.png} \\
\addlinespace
RF-01 & Periodo con categoría nueva &
Periodo nuevo y una categoría nueva escrita en el campo de categorías nuevas. &
Se crea la categoría y queda asignada al periodo. &
La categoría nueva se creó y quedó asignada al periodo. &
\ev{RF01\_periodo\_crear\_con\_categoria\_nueva.png} \\
\addlinespace
RF-01 & Fechas inválidas &
Fecha de fin anterior a la fecha de inicio. &
Mensaje de error y no se guarda. &
Mostró el mensaje de fecha de fin posterior a la de inicio y no guardó. &
\ev{RF01\_fechas\_error.png} \\
\addlinespace
RF-01 & Periodo repetido &
Nombre de un periodo que ya existe. &
Mensaje de error y no se guarda. &
Mostró ``Ya existe un periodo con ese nombre'' y no guardó. &
\ev{RF01\_periodo\_nombre\_repetido\_error.png} \\
\addlinespace
RF-01 & Editar periodo &
Cambio de nombre o de fechas. &
Se actualiza con el mensaje ``Periodo actualizado''. &
Se actualizó y apareció el mensaje. &
\ev{RF01\_periodo\_editar.png} \\
\addlinespace
RF-01 & Cerrar periodo &
Periodo en estado ACTIVO. &
El estado cambia a CERRADO. &
El estado cambió a CERRADO. &
\ev{RF01\_periodo\_cerrar.png} \\
\addlinespace
RF-01 & Eliminar periodo &
Periodo de prueba. &
Se elimina con el mensaje ``Periodo eliminado''. &
Se eliminó y apareció el mensaje. &
\ev{RF01\_periodo\_eliminar.png} \\
\addlinespace
RF-01 & Verificación en PostgreSQL &
Consulta con \texttt{JOIN} entre \texttt{periodo}, \texttt{periodo\_categoria} y \texttt{categoria\_documental}. &
Los datos de la base coinciden con lo hecho en la aplicación. &
Los registros coincidieron con lo hecho en la aplicación. &
\ev{RF01\_verificacion\_postgres\_periodo\_categoria.png} \\
\bottomrule
\end{longtable}
\endgroup